\documentclass[11pt,a4paper]{article}

\usepackage{iftex}

\ifpdftex
  \usepackage[utf8]{vietnam}
  \usepackage{tgpagella}
  \usepackage{amsmath, amssymb}
\else
  \usepackage{fontspec}
  \usepackage{unicode-math}
  \setmainfont{texgyrepagella-regular.otf}[
      BoldFont = texgyrepagella-bold.otf,
      ItalicFont = texgyrepagella-italic.otf,
      BoldItalicFont = texgyrepagella-bolditalic.otf
  ]
  \setmathfont{texgyrepagella-math.otf}
  \usepackage{amsmath}
\fi

\usepackage[a4paper, margin=2.5cm]{geometry}
\usepackage{microtype}
\usepackage{titlesec}
\usepackage{enumitem}

% Căn chỉnh khoảng cách theo chuẩn quy định
\linespread{1.05}
\setlength{\parindent}{0pt}
\setlength{\parskip}{5pt plus 2pt minus 1pt}

% Khoảng cách hiển thị công thức toán
\setlength{\abovedisplayskip}{5pt plus 2pt minus 2pt}
\setlength{\belowdisplayskip}{5pt plus 2pt minus 2pt}
\setlength{\abovedisplayshortskip}{3pt plus 1pt}
\setlength{\belowdisplayshortskip}{3pt plus 1pt}

% Định dạng tiêu đề câu hỏi
\titleformat{\section}
  {\large\bfseries}
  {\thesection}{0.5em}{}
\titlespacing*{\section}
  {0pt}{12pt plus 3pt minus 2pt}{4pt plus 1pt minus 1pt}

\pagestyle{plain}

\begin{document}

\begin{center}
    {\Large\bfseries ĐỀ THI VẬT LÝ ĐIỆN TỬ (PH3330)}\\[0.2cm]
    \textit{(Thời gian làm bài: 90 phút --- Không sử dụng tài liệu)}
\end{center}
\vspace{0.2cm}

%=====================================================================
% CÂU 1
%=====================================================================
\section*{Câu 1 (2 điểm). Dẫn sóng bản phẳng ba lớp}

So sánh sự lan truyền ánh sáng trong linh kiện dẫn sóng bản phẳng ba lớp bằng phương pháp quang học vật lý và quang học tia.

\vspace{0.5cm}

%=====================================================================
% CÂU 2
%=====================================================================
\section*{Câu 2 (2 điểm). Phổ năng lượng của phân tử}

Hãy trình bày tóm tắt về phổ năng lượng của phân tử.

\vspace{0.5cm}

%=====================================================================
% CÂU 3
%=====================================================================
\section*{Câu 3 (3 điểm). Nồng độ hạt dẫn trong silic pha nhôm}

Một mẫu tinh thể silic có khối lượng là $100\text{ g}$ được pha tạp $4\text{ }\mu\text{g}$ nhôm một cách đồng đều. Xác định nồng độ lỗ trống, điện tử, mức Fermi ở $300\text{ K}$. Cho biết khối lượng riêng của silic là $\rho_{\text{Si}} = 1{,}23\text{ g/cm}^3$, nguyên tử lượng của silic là $A_{\text{Si}} = 20{,}086$, nguyên tử lượng của nhôm là $A_{\text{Al}} = 26{,}98$.

\vspace{0.5cm}

%=====================================================================
% CÂU 4
%=====================================================================
\section*{Câu 4 (3 điểm). Hấp thụ ánh sáng và tái hợp trong GaAs}

Một mẫu GaAs được chiếu sáng bởi ánh sáng có bước sóng $\lambda = 0{,}5\text{ }\mu\text{m}$. Công suất tia tới là $25\text{ mW}$. Xác định chiều dày của mẫu nếu $1/4$ công suất đó bị phản xạ trở lại và $1/3$ thoát ra từ mặt sau của mẫu. Tìm năng lượng nhiệt được truyền cho mạng tinh thể trong $1\text{ s}$ và số photon sinh ra do tái hợp bức xạ riêng trong $100\text{ s}$. Giả thiết hệ số hấp thụ $\alpha(0{,}5\text{ }\mu\text{m}) = 4 \cdot 10^4\text{ cm}^{-1}$, $\Delta E_g = 1{,}42\text{ eV}$.

\end{document}
